% Sân tập SumiTeX. Đây là bản riêng của bạn, muốn sửa gì cũng được.
% (Your own copy of the SumiTeX playground: edit anything.)
%
% Phiên bản sân tập: 2026.11
%
% Nhờ dòng trên, SumiTeX biết khi nào có bài tập mới. Quản trị viên sẽ đổi
% số này mỗi lần cập nhật bài. (The server compares this line with the
% master playground to say "New exercises".)

\documentclass[11pt]{article}
\usepackage[T5]{fontenc}
\usepackage[utf8]{inputenc}
\usepackage[vietnamese]{babel}
% Để in được các ký hiệu giao diện trong bài tập (✎ ⠿ ⋯ ✓, mũi tên, α)
\usepackage{amssymb,pifont}
\DeclareUnicodeCharacter{270E}{\ding{46}}
\DeclareUnicodeCharacter{2713}{\ding{51}}
\DeclareUnicodeCharacter{283F}{\raisebox{-0.2ex}{$\vdots\!\vdots$}}
\DeclareUnicodeCharacter{22EF}{\ensuremath{\cdots}}
\DeclareUnicodeCharacter{2191}{\ensuremath{\uparrow}}
\DeclareUnicodeCharacter{2192}{\ensuremath{\rightarrow}}
\DeclareUnicodeCharacter{03B1}{\ensuremath{\alpha}}
\usepackage[a4paper,margin=2.2cm]{geometry}
\usepackage{graphicx}
\usepackage{amsmath}
% Mỗi việc cần làm là một \todo{…}. Trong PDF nó là một khung màu cam, còn
% trong Mục lục nó là một dòng ghi chú. Làm xong thì xoá đi, hoặc bấm
% ✎ rồi ✓ Done trong Mục lục.
\usepackage[colorinlistoftodos]{todonotes}
% Khung nằm ngay chỗ cần làm và rộng bằng trang, không nằm ngoài lề
\presetkeys{todonotes}{inline,color=orange!25,bordercolor=orange}{}

\title{Sân tập SumiTeX}
\author{Tên của bạn}
\date{}

\begin{document}
\maketitle
\listoftodos[Việc còn phải làm]

% Mỗi bài nằm trên một trang riêng
\newpage% Bài 1: tài liệu đầu tiên (xem phần "Văn bản đầu tiên" trong sổ tay)
\section*{Bài 1: Văn bản}

\todo{Sửa \textbackslash author trong main.tex thành tên của bạn}
\todo{Thêm mục \textbackslash section\{Giới thiệu\} và viết vài câu}
\todo{In đậm một từ bất kỳ bằng Ctrl+B}

% Gợi ý: một mục mới trông như thế này
%   \section{Giới thiệu}
%   Đây là đoạn văn đầu tiên.

Bạn viết vào đây nhé.

\newpage% Bài 2: danh sách (xem phần "Danh sách, bảng và hình" trong sổ tay)
\section*{Bài 2: Danh sách}

\todo{Thêm ý thứ ba vào danh sách bên dưới}
\todo{Chuyển sang danh sách đánh số (bấm ✎ trong Mục lục, chọn Numbered)}
\todo{Dùng Alt+↑ để đưa ý ``Làm bài'' lên đầu}

\begin{itemize}
  \item Mở sổ tay
  \item Làm bài
\end{itemize}

\newpage% Bài 3: bảng (xem phần "Danh sách, bảng và hình" trong sổ tay)
\section*{Bài 3: Bảng}

\todo{Thêm hàng ``Văn'' được 8 điểm (bấm ✎ ở bảng, rồi bấm + Row)}
\todo{Đặt chú thích là ``Điểm của tôi'' và nhãn là tab:diem}
\todo{Kéo dấu ⠿ để đổi chỗ hai cột}

\begin{tabular}{lc}
  Môn & Điểm \\
  Toán & 9 \\
\end{tabular}

\newpage% Bài 4: hình (xem phần "Danh sách, bảng và hình" trong sổ tay)
\section*{Bài 4: Hình}

\todo{Chụp một màn hình bất kỳ rồi nhấn Ctrl+V để dán vào dòng trống bên dưới}
\todo{Chỉnh hình rộng ½ bằng thanh nhỏ phía trên hình}
\todo{Gõ chú thích cho hình}

% Dán ảnh vào dòng trống dưới đây:


\newpage% Bài 5: công thức (xem phần "Công thức toán" trong sổ tay)
\section*{Bài 5: Công thức}

\todo{Viết định lý Pythagoras $a^2 + b^2 = c^2$ trong một equation có nhãn}
\todo{Thêm một phân số bằng hàng nút trong ô sửa công thức}
\todo{Chèn chữ α từ bảng ký hiệu (Ctrl+K rồi M)}
\todo{Dùng nút ⋯ rồi Insert a reference để dẫn tới công thức vừa viết}

% Gợi ý: một công thức có đánh số
%   \begin{equation}
%     ...
%     \label{eq:ten}
%   \end{equation}

Bạn viết vào đây nhé.


% Bài 6 cố tình có lỗi nên tạm tắt. Khi muốn làm, xoá dấu % ở đầu dòng dưới.
% \newpage% Bài 6: tự sửa lỗi (xem phần "Khi có lỗi" trong sổ tay)
%
% Bài này CỐ TÌNH có lỗi. Biên dịch (Ctrl+S), mở bảng lỗi (Ctrl+Shift+M)
% rồi sửa lần lượt, bắt đầu từ lỗi đầu tiên.
\section*{Bài 6: Sửa lỗi}

\todo{Sửa hết lỗi trong bài này cho tới khi bảng lỗi trống trơn}

Tập số thực là $\mathbb{R}$.

Điều này rất \textbf{quan trọng.

Bình phương của x là x^2.

\secton{Kết luận}

\todo{Bài 6: xoá dấu \% trước dòng \textbackslash input\{bai-06\} trong main.tex rồi sửa hết lỗi}

\end{document}
